% ID: FIS-MAG-FAR-002
% Titolo: Circuitazione media del campo elettrico indotto in una spira
% Disciplina: Fisica
% Area: Magnetismo
% Argomento: Faraday-Lenz
% Sottoargomento: Variazione del flusso magnetico in una spira
% Classe: Quinta liceo scientifico
% Difficolta: 4
% Tipo: problema_numerico
% Risultato: circa 7,6 * 10^{-9} V
% Tempo_stimato: 12 min
% Competenze: calcolo del flusso magnetico, applicazione della legge di Faraday, modellizzazione
% Prerequisiti: flusso del campo magnetico, legge di Faraday-Lenz, area del cerchio
% Tag: fisica, magnetismo, faraday_lenz, flusso_magnetico, spira
% Fonte: verifica_fisica_magnetismo_anonimizzata
% Autore: Riccardo Carli
% Licenza: CC-BY-SA-4.0

\begin{esercizio}
Una spira circolare di raggio \(8{,}0\,\mathrm{cm}\) e immersa in un campo
magnetico uniforme di intensita \(B = 3{,}8 \cdot 10^{-6}\,\mathrm{T}\).
All'istante \(t=0\,\mathrm{s}\) la spira e perpendicolare al campo magnetico;
viene poi fatta ruotare lentamente intorno a un suo diametro fino a essere
inclinata di \(60^\circ\) rispetto alla direzione iniziale all'istante
\(t=5{,}0\,\mathrm{s}\).

Calcola il modulo della circuitazione media del campo elettrico lungo la spira.
\end{esercizio}

\begin{soluzione}
La circuitazione del campo elettrico indotto e uguale alla forza elettromotrice
indotta:
\[
\left|\oint \vec E\cdot d\vec l\right|
=\left|\frac{\Delta\Phi_B}{\Delta t}\right|.
\]

L'area della spira e:
\[
A=\pi r^2=\pi(0{,}080)^2
\simeq 2{,}01\cdot 10^{-2}\,\mathrm{m^2}.
\]

All'inizio la spira e perpendicolare al campo magnetico, quindi il suo vettore
superficie e parallelo a \(\vec B\) e il flusso e:
\[
\Phi_i=BA.
\]
Dopo la rotazione di \(60^\circ\), il flusso diventa:
\[
\Phi_f=BA\cos 60^\circ=\frac{1}{2}BA.
\]

Il modulo della variazione di flusso e:
\[
|\Delta\Phi_B|=\left|\Phi_f-\Phi_i\right|
=\frac{1}{2}BA.
\]

Quindi:
\[
\left|\oint \vec E\cdot d\vec l\right|
=\frac{\frac{1}{2}BA}{\Delta t}
=\frac{0{,}5\cdot 3{,}8\cdot 10^{-6}\cdot
2{,}01\cdot 10^{-2}}{5{,}0}.
\]

\[
\left|\oint \vec E\cdot d\vec l\right|
\simeq 7{,}6\cdot 10^{-9}\,\mathrm{V}.
\]
\end{soluzione}
